This work revisits the undergraduate thesis of W.~L.~S. de Paula (ITA,
2003), which found six polarization states for gravitational waves in
Visser's bimetric massive-graviton theory, and carries its results to
the observable of pulsar timing arrays (PTAs). Redone without the
quasi-null approximation, the analysis confirms exactly six independent
tidal amplitudes. It also shows that the sixth is the trace of the
perturbation, a scalar ghost that allows the model to evade the van
Dam--Veltman--Zakharov discontinuity. In the consistent Fierz--Pauli
theory, five polarizations remain, and the transverse and longitudinal
scalar deformations are tied by $p_l=-(f_g/f)^2p_b$. The Newman--Penrose
amplitudes of the thesis are exact contractions of the Riemann tensor,
but their reading as polarization amplitudes and the E(2)
classification assume quasi-null propagation and apply only for
$f\gg f_g$. For the low-frequency example of the 2004 paper, this
reading is off by about 60\%.
The $(f_g/f)^2$ suppression of the extra modes found in the thesis is
exact for comparable metric amplitudes. With current mass bounds,
however, specifically massive signatures are only appreciable in the
nanohertz band. The PTA response is shown to depend on the wave only
through the tidal matrix computed in the thesis. The angular
correlations of the tensor and helicity-zero sectors are computed with
pulsar terms, and a symmetry argument explains why, at the threshold
$f=f_g$, all polarizations yield the same correlation per state.
Controlled simulations show that small arrays learn little about the
mass (0.07 to 0.11 nat), that a fixed average of the correlations over frequencies removes
about 95\% of this information, and that the usual normal approximation for
correlation estimators increases the rate of false scalar-polarization
detections tenfold.
